\section*{Hoofdstuk \ref{ch:g8:fronts-and-weather-maps} --- Luchtmassa's, fronten en weerkaarten}

\begin{solution}{exo:g8:fronts-and-weather-maps:1}
(a) Maritiem tropisch: warm en vochtig. (b) Continentaal polair: koud en
droog. (c) Continentaal tropisch: heet en droog. (d) Maritiem polair: koud
en vochtig.
\end{solution}

\begin{solution}{exo:g8:fronts-and-weather-maps:2}
Blauwe lijn met driehoekjes: een \omterm{def:g8:fronts-and-weather-maps:front}{koufront}. Rode lijn met halve cirkels: een
\omterm{def:g8:fronts-and-weather-maps:front}{warmtefront}. Lijn die plaatsen met dezelfde druk verbindt: een \omterm{def:g8:fronts-and-weather-maps:isobar}{isobaar}. H:
het centrum van een \omterm{def:g8:fronts-and-weather-maps:depression}{anticycloon} (hoge druk). L: het centrum van een
\omterm{def:g8:fronts-and-weather-maps:depression}{depressie} (lage druk).
\end{solution}

\begin{solution}{exo:g8:fronts-and-weather-maps:3}
(a) Waar. (b) Niet waar: ze komen ervoor, ver vooruit, hoog in de lucht.
(c) Niet waar: \omterm{def:g8:fronts-and-weather-maps:isobar}{isobaren} dicht bij elkaar betekenen een sterke wind.
(d) Waar.
\end{solution}

\begin{solution}{exo:g8:fronts-and-weather-maps:4}
(a) \qty{980}{hPa}. (b) \qty{1034}{hPa}. (c) \qty{4}{hPa}. (d) $1034 -
980 = \qty{54}{hPa}$.
\end{solution}

\begin{solution}{exo:g8:fronts-and-weather-maps:5}
Dunne strepen \omterm{def:g4:clouds-rain-snow:cloud-families}{cirrus}; een grijze sluier van altostratus; de \omterm{def:g4:clouds-rain-snow:cloud}{wolken} komen
lager en worden dikker; gestage \omterm{def:g4:clouds-rain-snow:precipitation}{regen}; warmere, vochtigere lucht.
\end{solution}

\begin{solution}{exo:g8:fronts-and-weather-maps:6}
Bij een \omterm{def:g8:fronts-and-weather-maps:front}{koufront} schuift de dichte koude lucht als een \emph{wig} onder de
warme lucht: de warme lucht wordt langs een \emph{steile} helling
opgetild, stijgt snel op en bouwt hoge \omterm{def:g4:clouds-rain-snow:cloud}{wolken}. Bij een \omterm{def:g8:fronts-and-weather-maps:front}{warmtefront} klimt de
warme lucht langzaam over de koude lucht, langs een \emph{flauwe helling}:
ze stijgt beetje bij beetje op over een lange afstand, en haar \omterm{def:g4:clouds-rain-snow:cloud}{wolken}
spreiden zich in lagen uit.
\end{solution}

\begin{solution}{exo:g8:fronts-and-weather-maps:7}
(a) Bij L: daar liggen de \omterm{def:g8:fronts-and-weather-maps:isobar}{isobaren} dichter bij elkaar, en de pijlen zijn
langer. (b) Uit het oosten (de pijlen wijzen naar het westen, een beetje
naar het zuidwesten). (c) Uit het oosten (de pijlen wijzen naar het
westen). (d) Tegen de klok in.
\end{solution}

\begin{solution}{exo:g8:fronts-and-weather-maps:8}
(a) $1016 - 1003 = \qty{13}{hPa}$. (b) Een \omterm{def:g8:fronts-and-weather-maps:depression}{depressie}: de druk blijft \omterm{def:g1:landscapes-around-us:landform}{dalen}.
(c) Bewolking, wind en \omterm{def:g4:clouds-rain-snow:precipitation}{regen}, waarschijnlijk met een \omterm{def:g8:fronts-and-weather-maps:front}{front}.
\end{solution}

\begin{solution}{exo:g8:fronts-and-weather-maps:9}
(a) cT, continentaal tropisch. (b) Heet en droog. (c) Boven de \omterm{def:g2:water-on-earth:ocean}{zee} neemt
ze waterdamp op: ze wordt vochtiger, en de lucht dicht bij het water wordt
door de \omterm{def:g2:water-on-earth:ocean}{zee} een beetje afgekoeld.
\end{solution}

\begin{solution}{exo:g8:fronts-and-weather-maps:10}
(a) $4 - 1 = 3$ dagen. (b) $3 \div 3 = 1$ dag per tien jaar. (c) Er zijn
$5 - 1 = 4$ dagen meer nodig: ongeveer 40 jaar.
\end{solution}

\begin{solution}{exo:g8:fronts-and-weather-maps:11}
(a) Ongeveer \qty{1350}{km} (vierenhalve schaalbalk). (b) \qty{4}{hPa}
over \qty{80}{km}: $4 \div 80 \times 100 = \qty{5}{hPa}$ per
\qty{100}{km}. (c) $4 \div 140 \times 100 \approx \qty{2.9}{hPa}$ per
\qty{100}{km}. (d) Bij L, waar de druk sneller verandert; ja, de pijlen
zijn bij L langer.
\end{solution}

\begin{solution}{exo:g8:fronts-and-weather-maps:12}
(a) Het \omterm{def:g8:fronts-and-weather-maps:front}{koufront}: $300 \div 50 = 6$ uur. Het \omterm{def:g8:fronts-and-weather-maps:front}{warmtefront} verplaatst zich
met \qty{25}{km} per uur: $200 \div 25 = 8$ uur. Het \omterm{def:g8:fronts-and-weather-maps:front}{koufront} komt 2 uur
eerder aan. (b) Eerste stad: hoog oprijzende \omterm{def:g4:clouds-rain-snow:cloud}{wolken}, hevige buien,
windstoten, daarna koelere lucht en opklaringen. Tweede stad:
steeds lagere gelaagde \omterm{def:g4:clouds-rain-snow:cloud}{wolken} en een lange, gestage \omterm{def:g4:clouds-rain-snow:precipitation}{regen}, daarna warmere,
vochtigere lucht. (c) De tweede stad, die door warme lucht wordt bereikt.
\end{solution}

\begin{solution}{pb:g8:fronts-and-weather-maps:1}
\textbf{1.} In de warme lucht, tussen de twee \omterm{def:g8:fronts-and-weather-maps:front}{fronten}. Het \omterm{def:g8:fronts-and-weather-maps:front}{warmtefront} is
er al overheen getrokken.

\textbf{2.} Uit het zuidwesten.

\textbf{3.} Naar het oosten: de driehoekjes van een \omterm{def:g8:fronts-and-weather-maps:front}{koufront} staan aan de
kant waarheen het zich verplaatst, en ze wijzen naar het oosten, naar de
stad.

\textbf{4.} Hoog oprijzende \omterm{def:g4:clouds-rain-snow:cloud}{wolken}, hevige buien, misschien onweer en
\omterm{def:g4:clouds-rain-snow:precipitation}{hagel}, windstoten; daarna koelere lucht, een wind uit het westen of
noordwesten en opklaringen.

\textbf{5.} Op \qty{180}{km} om 6.00~uur en op \qty{300}{km} om 9.00~uur.

\textbf{6.} $300 - 180 = \qty{120}{km}$.

\textbf{7.} $120 \div 3 = \qty{40}{km}$ per uur.

\textbf{8.} De stad ligt op \qty{480}{km}: $480 - 300 = \qty{180}{km}$.

\textbf{9.} $180 \div 40 = 4.5$ uur, dat is 4 uur en 30 minuten.

\textbf{10.} 9.00~uur + 4~uur~30~min = \textbf{13.30~uur}.

\textbf{11.} Nee: de buien zouden rond 14.30~uur moeten ophouden, voor de
aftrap.

\textbf{12.} $180 \div 30 = 6$ uur: om 15.00~uur, precies bij de aftrap.
Verwachting: ``Het \omterm{def:g8:fronts-and-weather-maps:front}{koufront} zal de stad rond \textbf{13.30~uur} bereiken,
met een uur hevige buien en windstoten, daarna klaart het op: de wedstrijd
kan zeer waarschijnlijk worden gespeeld. Als het \omterm{def:g8:fronts-and-weather-maps:front}{front} vertraagt, kan het
om 15.00~uur aankomen en de wedstrijd vertragen.''
\end{solution}
